\documentclass[11pt,a4paper]{report}

\usepackage[T1]{fontenc}
\usepackage[english,french]{babel}
\usepackage[left=3cm,right=3cm]{geometry}
\usepackage{amsmath} 
\usepackage{amsfonts}
\usepackage{graphicx} 
\usepackage{setspace}

%%%%%%%%%%%%%%%%%%%%%%%%%%%%%%%%%%%%%%%%%%%%%%%%%%%%%%%%%%%%%%%%%%%%%%%%%%%%%%%%
  
\begin{document} 

\begin{titlepage}
\begin{center}
{\huge\bfseries Titre du compte-rendu} \\
% ----------------------------------------------------------------
\vspace{1.5cm}
{\Large\bfseries Prénom Nom}\\[5pt]
adresse.email@unistra.fr\\[14pt]
% ----------------------------------------------------------------
\vfill
UFR de Physique et Ingénierie \\[14pt]
Université de Strasbourg\\[14pt]
Année académique 2022-23
\end{center}
\end{titlepage}

%%%%%%%%%%%%%%%%%%%%%%%%%%%%%%%%%%%%%%%%%%%%%%%%%%%%%%%%%%%%%%%%%%%%%%%%%%%%%%%%

\newpage
\onehalfspacing

\begin{abstract}
  Résume (maximum 1/2 page) du travail effectue (contexte, résultats).
\end{abstract}

%%%%%%%%%%%%%%%%%%%%%%%%%%%%%%%%%%%%%%%%%%%%%%%%%%%%%%%%%%%%%%%%%%%%%%%%%%%%%%%%

\setcounter{page}{3}
\tableofcontents

%%%%%%%%%%%%%%%%%%%%%%%%%%%%%%%%%%%%%%%%%%%%%%%%%%%%%%%%%%%%%%%%%%%%%%%%%%%%%%%%

\chapter*{Introduction}
\addcontentsline{toc}{chapter}{Introduction}

L'introduction (maximum 2 pages) sert à donner le cadre du travail
(domaine de recherche dans lequel il se situe) et l'état de l'art. Il est normal de citer les références associées: articles\cite{Article01}, livres\cite{Livre01}, chapitres de livres\cite{Chapitre01}, conférences\cite{Proceeding01} ou thèses\cite{These01}.
Elle spécifie aussi comment les chapitres suivants sont structurés.
Tout comme le chapitre final ``Conclusions'', l'Introduction est
rédigée après tout le reste.


%%%%%%%%%%%%%%%%%%%%%%%%%%%%%%%%%%%%%%%%%%%%%%%%%%%%%%%%%%%%%%%%%%%%%%%%%%%%%%%%

\chapter{Bases théoriques / Connaissances de base / ....}

La première partie du texte a comme objectif de fournir un rappel des
connaissances nécessaires pour comprendre le travail effectue.  Le
titre tout comme son contenu dépendent du type de travail effectue.
Dans le cas d'un travail théorique ou de simulation, il est important de fournir les équations qui seront utilises dans les calculs ou simulations.  Si on s'appuie sur des travaux précédents, il faut (logiquement) les citer\cite{Article01}.  Pour un travail expérimental, il est important de
rappeler ici les phénomènes physiques qui seront utilisés.  Dans le
cas ou on s'appuie sur un système expérimental déjà existant, on
utilise ce chapitre pour en donner une description.


%%%%%%%%%%%%%%%%%%%%%%%%%%%%%%%%%%%%%%%%%%%%%%%%%%%%%%%%%%%%%%%%%%%%%%%%%%%%%%%%

\chapter{Algorithmes / Système expérimental / ....}

Le cadre théorique défini, on passe a une description plus précise du
travail spécifique effectué.  Tout comme dans les sections précédentes, il est normal de faire référence\cite{Livre01,Chapitre01} aux textes d'où on a pris les méthodes utilisées.

\begin{itemize}
\item Dans le cas d'un travail théorique on montre ici les calculs
  effectues et les étapes pour traiter / simplifier / résoudre les
  équations.
\item Pour des simulations, on discute des algorithmes numériques, 
  des choix effectues, approximations, etc.
\item Pour un travail expérimental on présente la mise en place du
  système (si elle fait partie du travail), les procédures de mesure
  et de traitement des données.
\end{itemize}

Le chapitre sert a montrer la démarche utilisée pour arriver aux
résultats qui seront ensuite analyses et discutes dans le chapitre
suivant.


%%%%%%%%%%%%%%%%%%%%%%%%%%%%%%%%%%%%%%%%%%%%%%%%%%%%%%%%%%%%%%%%%%%%%%%%%%%%%%%%

\chapter{Résultats / Analyse}

Dans ce dernier chapitre on fournit les résultats obtenus et on les
analyse afin de mettre en évidence les nouvelles connaissances
générées et acquises lors du stage.


%%%%%%%%%%%%%%%%%%%%%%%%%%%%%%%%%%%%%%%%%%%%%%%%%%%%%%%%%%%%%%%%%%%%%%%%%%%%%%%%

\chapter*{Conclusions}
\addcontentsline{toc}{chapter}{Conclusions}

Les conclusions sont le symétrique de l'Introducion: on résume ici le
chemin effectue, les résultats marquants et les éventuels
développements futurs que ces résultats vont permettre.

%%%%%%%%%%%%%%%%%%%%%%%%%%%%%%%%%%%%%%%%%%%%%%%%%%%%%%%%%%%%%%%%%%%%%%%%%%%%%%%%

\clearpage
\addcontentsline{toc}{chapter}{Bibliographie}
\bibliography{biblio}
\bibliographystyle{unsrt}

Le travail scientifique n'existe pas dans le vide et s'appuie toujours
sur des connaissances préexistantes.  On liste ici les textes (livres,
articles, conférences) auxquels on a fait référence dans le texte.


%%%%%%%%%%%%%%%%%%%%%%%%%%%%%%%%%%%%%%%%%%%%%%%%%%%%%%%%%%%%%%%%%%%%%%%%%%%%%%%%

\appendix
\chapter{Détails additionnels}

Les annexes (une ou plusieurs) servent à indiquer des détails qui
sont importants pour reproduire le travail mais qui alourdiraient trop
le texte principal.

\end{document}
